\begin{thebibliography}{99}
\addcontentsline{toc}{chapter}{Bibliografía}

% --- Requisitos y competencias
\bibitem{anexo} Anexo de requisitos técnicos, CanSat Universitario 2026--2027.
  Documento de la organización entregado a los equipos.

\bibitem{cansat2019} CanSat Competition. \emph{CanSat Competition Guide 2019. Mission:
  Auto-gyro Probe}, rev.~1.1a, 21 de julio de 2018. American Astronautical Society y
  otros. \url{https://www.readkong.com/page/cansat-competition-guide-2019-1144438}

\bibitem{cansat2025} CanSat Competition. \emph{CanSat Competition Guide 2025. Auto Gyro
  Descender}, rev.~1.0d, 31 de octubre de 2024.
  \url{https://cansatcompetition.com/docs/CanSat_Mission_Guide_2025d.pdf}

% --- Artículos sobre CanSat con autogiro
\bibitem{ventura2022} A.~Ventura, J.~Nuñez-Quispe, G.~Sanchez y J.~Santivañez.
  CanSat payload with Autogyro for descent in experimental rocket flights: Development,
  CFD analysis, and preliminary test on free-fall. \emph{Journal of Physics: Conference
  Series}, 2235:012055, 2022. doi:10.1088/1742-6596/2235/1/012055.

\bibitem{ramadhan2019} R.~P.~Ramadhan, A.~R.~Ramadhan, S.~A.~Putri y otros. Prototype of
  CanSat with Auto-gyro Payload for Small Satellite Education. En \emph{2019 IEEE 13th
  International Conference on Telecommunication Systems, Services, and Applications
  (TSSA)}, pp.~243--248, 2019. doi:10.1109/TSSA48701.2019.8985514.

\bibitem{chun2023} C.~Chun, M.~H.~Tanveer y S.~Chakravarty. The CanSat Compendium: A
  Review of Scientific CanSats. \emph{Machines}, 11(7):675, 2023.

% --- Aerodinámica de rotores
\bibitem{glauert1926} H.~Glauert. A General Theory of the Autogyro. \emph{Aeronautical
  Research Committee, Reports and Memoranda} No.~1111, 1926.

\bibitem{castles1951} W.~Castles Jr. y R.~B.~Gray. Empirical Relation Between Induced
  Velocity, Thrust and Rate of Descent of a Helicopter Rotor as Determined by Wind Tunnel
  Tests on Four Model Rotors. \emph{NACA TN~2474}, 1951.

\bibitem{gessow1952} A.~Gessow y G.~C.~Myers. \emph{Aerodynamics of the Helicopter}.
  Macmillan, 1952 (reimpr. Frederick Ungar, 1967).

\bibitem{johnson1980} W.~Johnson. \emph{Helicopter Theory}. Princeton University Press,
  1980 (reimpr. Dover, 1994).

\bibitem{johnson2013} W.~Johnson. \emph{Rotorcraft Aeromechanics}. Cambridge University
  Press, 2013.

\bibitem{leishman2006} J.~G.~Leishman. \emph{Principles of Helicopter Aerodynamics},
  2.ª ed. Cambridge University Press, 2006.

\bibitem{leishman2004} J.~G.~Leishman, M.~J.~Bhagwat y S.~Ananthan. The Vortex Ring
  State as a Spatially and Temporally Developing Wake Instability. \emph{Journal of the
  American Helicopter Society}, 49(2):160--175, 2004.

\bibitem{johnson2005} W.~Johnson. Model for Vortex Ring State Influence on Rotorcraft
  Flight Dynamics. \emph{NASA/TP-2005-213477}, NASA Ames Research Center, 2005.

% --- Aerodinámica de bajo número de Reynolds
\bibitem{lissaman1983} P.~B.~S.~Lissaman. Low-Reynolds-Number Airfoils. \emph{Annual Review
  of Fluid Mechanics}, 15:223--239, 1983.

\bibitem{mueller2003} T.~J.~Mueller y J.~D.~DeLaurier. Aerodynamics of Small Vehicles.
  \emph{Annual Review of Fluid Mechanics}, 35:89--111, 2003.

\bibitem{pelletier2000} A.~Pelletier y T.~J.~Mueller. Low Reynolds Number Aerodynamics of
  Low-Aspect-Ratio, Thin/Flat/Cambered-Plate Wings. \emph{Journal of Aircraft},
  37(5):825--832, 2000.

\bibitem{okamoto1996} M.~Okamoto, K.~Yasuda y A.~Azuma. Aerodynamic Characteristics of the
  Wings and Body of a Dragonfly. \emph{Journal of Experimental Biology}, 199:281--294, 1996.

\bibitem{selig1995} M.~S.~Selig, J.~J.~Guglielmo, A.~P.~Broeren y P.~Giguère.
  \emph{Summary of Low-Speed Airfoil Data}, vol.~1. SoarTech Publications, 1995.

\bibitem{schmitz1942} F.~W.~Schmitz. \emph{Aerodynamik des Flugmodells}. C.~J.~E.~Volckmann,
  Berlín, 1942.

% --- Semillas autorrotantes
\bibitem{norberg1973} R.~Å.~Norberg. Autorotation, Self-Stability, and Structure of
  Single-Winged Fruits and Seeds (Samaras) with Comparative Remarks on Animal Flight.
  \emph{Biological Reviews}, 48(4):561--596, 1973.

\bibitem{azuma1989} A.~Azuma y K.~Yasuda. Flight Performance of Rotary Seeds.
  \emph{Journal of Theoretical Biology}, 138:23--53, 1989.

\bibitem{seter1992} D.~Seter y A.~Rosen. Study of the Vertical Autorotation of a
  Single-Winged Samara. \emph{Biological Reviews}, 1992.

\bibitem{lentink2009} D.~Lentink, W.~B.~Dickson, J.~L.~van~Leeuwen y M.~H.~Dickinson.
  Leading-Edge Vortices Elevate Lift of Autorotating Plant Seeds. \emph{Science},
  324(5933):1438--1440, 2009. doi:10.1126/science.1174196.

% --- Paracaídas, resistencia y estructura
\bibitem{knacke1992} T.~W.~Knacke. \emph{Parachute Recovery Systems Design Manual}.
  NWC TP~6575. Para Publishing, Santa Bárbara, 1992.

\bibitem{hoerner1965} S.~F.~Hoerner. \emph{Fluid-Dynamic Drag}. Publicado por el autor,
  1965.

\bibitem{timoshenko1961} S.~P.~Timoshenko y J.~M.~Gere. \emph{Theory of Elastic
  Stability}, 2.ª ed. McGraw-Hill, 1961.

\bibitem{budynas2015} R.~G.~Budynas y J.~K.~Nisbett. \emph{Shigley's Mechanical
  Engineering Design}, 10.ª ed. McGraw-Hill, 2015.

\bibitem{anderson2017} J.~D.~Anderson. \emph{Fundamentals of Aerodynamics}, 6.ª ed.
  McGraw-Hill, 2017.

\end{thebibliography}

\noindent\small\textit{Nota sobre las fuentes.} Las referencias de competencias y
artículos sobre CanSat se verificaron en línea el 2 de octubre de 2026. Las de libros y
artículos clásicos se citan por su edición conocida; antes de una entrega formal conviene
confirmar páginas y ediciones contra los ejemplares disponibles.
